% ECONOMICS_PROBLEM: {"id": "L41-501", "title": "Merger efficiencies", "jel_code": "L41", "source_id": "OP-0394"}
% FIRST_POSTED: 2026-10-09
% ATTEMPT_STATUS: partial
% ENTRY_KIND: research_attempt
\documentclass[11pt,a4paper]{article}
\usepackage[T1]{fontenc}
\usepackage[utf8]{inputenc}
\usepackage[margin=27mm]{geometry}
\usepackage{amsmath,amssymb,amsthm,mathtools}
\usepackage[hidelinks]{hyperref}
\setlength{\parindent}{0pt}
\setlength{\parskip}{5pt}
\newtheorem{theorem}{Theorem}[section]
\newtheorem{proposition}[theorem]{Proposition}
\newtheorem{lemma}[theorem]{Lemma}
\newcommand{\R}{\mathbb R}
\newcommand{\E}{\mathbb E}
\DeclareMathOperator{\Var}{Var}
\DeclareMathOperator{\Cov}{Cov}
\begin{document}
\section*{Merger efficiencies: calibration, coupling, and a price threshold}
\section{What is established}
The catalogue asks about a population of proposed deals and annual real costs at a common reference output. Neither profit growth nor a change in average accounting cost is that outcome. This attempt gives an exact finite-support characterization of the joint savings and consumer-benefit distribution, a conditional calibration result, and a pricing example showing why positive efficiencies need not benefit consumers. The empirical classification of a five-year deal cohort remains unresolved. No merger audit observations are supplied here, and the numerical examples below are constructed examples.

Write $Z_a=(C_a,P_a)$ for the cost and fixed-basket price vector under regime $a$. Fix a preapproval covariate cell $(s,x)$ with positive population probability. Assume consistency, conditional independence $(Z_0,Z_1)\perp A\mid S,X$, and $0<\Pr(A=1\mid s,x)<1$. These are substantive assumptions about approval or abandonment; a regulator's use of private efficiency evidence can violate them. Assume all means exist. Then the observable conditional laws identify the two marginal distributions $\mu_a$ of $Z_a$, and
\[
 g(s,x)=E[C\mid A=0,s,x]-E[C\mid A=1,s,x]. \tag{1}
\]
The proof is iterated expectation under independence and consistency. Calibration error equals (1) minus $s$, provided the claimed savings use exactly the same annual resource-cost and output convention. This argument identifies a difference of means without identifying which potential cost outcomes belong to the same deal. It does not license interpreting an estimated conditional function off the observed support.

\section{An exact transport description of joint benefits}
Suppose, for a finite benchmark, $\mu_0$ puts masses $u_i$ on pairs $(c_{0i},p_{0i})$, $i\le I$, and $\mu_1$ puts masses $v_j$ on $(c_{1j},p_{1j})$, $j\le J$. Let
\[
 \mathcal T(u,v)=\{T\ge0:\ \sum_jT_{ij}=u_i,\quad\sum_iT_{ij}=v_j\}.
\]
This is a nonempty compact polytope: $T_{ij}=u_iv_j$ is feasible. Every $T$ induces a distribution of $(G,B)$ by assigning mass $T_{ij}$ to
$(c_{0i}-c_{1j},p_{0i}-p_{1j})$.
\begin{theorem}
Under the preceding assumptions and no cross-regime restrictions, the set of these pushforward distributions is exactly the joint identified set. In particular, for any specified function $h$ on the finite difference support, sharp lower and upper bounds on $E[h(G,B)]$ are the minimum and maximum of
\[
 \sum_{i,j}T_{ij}h(c_{0i}-c_{1j},p_{0i}-p_{1j}) \tag{2}
\]
over $\mathcal T(u,v)$.
\end{theorem}
\begin{proof}
Any compatible potential-outcome law supplies its joint mass table $T$, whose margins must be $u,v$. Conversely, draw a pair $(Z_0,Z_1)$ from any feasible table, independently assign $A$ with the stipulated cell propensity, and report $Z_A$. This model has precisely the observed law and the induced distribution of differences. Compactness gives attainment of (2). Randomizing between two feasible tables preserves both margins, so every intermediate value of (2) is attainable as well.
\end{proof}
Taking $h=\mathbf1\{G>0,B>0\}$ answers the probability of simultaneous benefits; taking $h=GB$ bounds their product moment. The same table must be used for all functionals. Separate scalar bounds generally do not form their joint feasible set. Rank invariance, engineering restrictions, or a common latent-type model can be imposed as additional constraints, but their economic content must be supplied before observing the desired conclusion. A zero restriction $T_{ij}=0$ is a transparent way to encode a genuinely impossible cross-regime pairing.

For example, let both regime distributions assign probability one-half to $(C,P)=(0,0)$ and one-half to $(2,2)$. Pairing like points gives $(G,B)=(0,0)$ surely. Pairing opposite points gives $(G,B)=(2,2)$ and $(-2,-2)$ with equal probability. Both regimes and both average effects are identical in the two models, yet the probability of strict simultaneous benefit is zero versus one-half. In fact one-half is the sharp upper bound: simultaneous positivity requires precisely the row $(2,2)$ and column $(0,0)$, each of mass one-half. This example preserves within-regime cost-price dependence, which would be lost by coupling four univariate margins separately.

\section{Selection and unavailable comparisons}
If the independence condition is unjustified, observed approval groups do not identify $\mu_a$ for the full reviewed population. For a bounded scalar outcome $C_a\in[\ell,u]$, let $\pi=\Pr(A=1)$ and $m_a=E[C\mid A=a]$. Without selection restrictions,
\[
 E[C_0]\in[(1-\pi)m_0+\pi\ell,(1-\pi)m_0+\pi u],
\]
\[
 E[C_1]\in[\pi m_1+(1-\pi)\ell,\pi m_1+(1-\pi)u].
\]
Subtracting opposite endpoints gives sharp bounds on average savings: assign all missing potential costs to the corresponding endpoint, leaving the observed coordinates unchanged. Price outcomes can be bounded analogously, but joint savings-price restrictions again require a common latent table. If no credible finite resource-cost bound exists, finite numerical bounds cannot be produced by the formula alone. An audit of selected successful mergers does not repair selection into approval or into the audit.

For verified fixed covariate bins, independent observations and $C\in[\ell,u]$ yield a simple simultaneous calibration interval. With $K$ bins and $n_{ak}>0$ observations in each regime-bin cell, define
\[
 r_{ak}=(u-\ell)\sqrt{\log(4K/\alpha)/(2n_{ak})}.
\]
Hoeffding and a union bound give coverage at least $1-\alpha$ for all conditional savings means using radii $r_{0k}+r_{1k}$ around the differences of sample means. Empty cells receive their full feasible bounds. The assertion concerns bin averages under independent representative sampling and valid selection assumptions; it is not a confidence band for an unrestricted continuous regression surface.

\section{Efficiencies sufficient to offset market power}
Consider inverse demand $P=a-Q$, two premerger Cournot firms with constant marginal cost $c<a$, and no capacity constraint. Each supplies $(a-c)/3$ and the premerger price is $(a+2c)/3$. The merged monopolist has marginal cost $c-e$, where $e\ge0$ and $c-e\ge0$, and chooses output $(a-c+e)/2$. Thus
\[
 B=P_0-P_1=\frac{3e-(a-c)}6. \tag{3}
\]
At reference output $\bar Q>0$, technological variable-cost savings are $e\bar Q$. A reduction $f$ in annual fixed resource cost increases $G$ to $e\bar Q+f$, but does not enter (3). These equations follow directly from the firms' first-order conditions, with strict concavity guaranteeing the respective optima. They exhibit positive fixed-output savings and higher consumer prices whenever $G>0$ but $e<(a-c)/3$. Claimed fixed-cost reductions therefore cannot simply be converted into marginal-cost reductions in a pricing simulation. The benchmark is deliberately a two-firm-to-monopoly comparison; it supplies neither a universal pass-through coefficient nor an empirical merger verdict.

\section{Remaining research task}
A completed application needs deal-level production and logistics audits, cost bounds for missing regimes, a defensible approval comparison, and a coupling sensitivity analysis. Purchasing discounts caused by bargaining power must be separated from lower resource use; otherwise (1) estimates a different cost concept. The transport program solves the finite statistical problem once valid margins and restrictions are supplied, while the cohort's actual calibration and welfare remain unknown.
\begin{thebibliography}{9}
\bibitem{sy} C. Shapiro and A. Yurukoglu (2024), \emph{Trends in Competition in the United States: What Does the Evidence Show?}, author draft dated 24 August, Sections 4.C and 6. \url{https://web.stanford.edu/~ayurukog/trendsincompetition.pdf}. Source motivation for evaluating merger efficiencies; the transport and Cournot calculations above are explicit benchmarks, not estimates from that paper.
\end{thebibliography}

\end{document}
